% Part: sets-functions-relations
% Chapter: sets
% Section: russells-paradox

\documentclass[../../../include/open-logic-section]{subfiles}


\begin{document}

\olfileid[id]{sfr}{set}{rus}
\olsection{Paradoks Russell}

Ekstensionalitas membenarkan notasi $\Setabs{x}{\phi(x)}$ untuk
\emph{himpunan} semua $x$ sedemikian sehingga~$\phi(x)$. Akan tetapi,
ekstensionalitas \emph{sesungguhnya} hanya membenarkan pemikiran berikut.
\emph{Jika} ada himpunan yang anggotanya tepat semua objek yang
memenuhi $\phi$, \emph{maka} hanya ada satu himpunan semacam itu.
Dengan kata lain, setelah suatu~$\phi$ ditetapkan, himpunan
$\Setabs{x}{\phi(x)}$ bersifat unik \emph{jika himpunan itu ada}.

Namun, syarat ini penting!{} Tidak setiap sifat memungkinkan
\emph{komprehensi}, yaitu pembentukan himpunan dari semua objek yang
memiliki sifat tersebut. Dengan kata lain, beberapa sifat
\emph{tidak} mendefinisikan himpunan. Jika semuanya mendefinisikan
himpunan, kita akan menghadapi kontradiksi langsung. Contoh paling
terkenal adalah Paradoks Russell.

Himpunan dapat menjadi !!{element}s dari himpunan lain---sebagai
contoh, himpunan kuasa dari himpunan~$A$ tersusun atas himpunan-himpunan.
Karena itu, masuk akal untuk menanyakan atau menyelidiki apakah sebuah
himpunan merupakan !!a{element} dari himpunan lain. Dapatkah sebuah
himpunan menjadi anggota dirinya sendiri? Tidak ada sesuatu dalam
gagasan tentang himpunan yang tampaknya menutup kemungkinan ini.
Sebagai contoh, jika \emph{semua} himpunan membentuk suatu kumpulan
objek, mungkin kita berpikir bahwa semuanya dapat dihimpun menjadi
satu himpunan---himpunan semua himpunan. Karena himpunan itu sendiri
adalah sebuah himpunan, ia akan menjadi !!a{element} dari himpunan
semua himpunan.

Paradoks Russell muncul ketika kita meninjau sifat tidak memuat dirinya
sendiri sebagai !!a{element}, yakni sifat \emph{tidak beranggotakan
diri sendiri}. Bagaimana jika kita mengandaikan bahwa terdapat
himpunan semua himpunan yang tidak memuat dirinya sendiri sebagai
!!a{element}? Apakah
\[
R = \Setabs{x}{x \notin x}
\]
ada? Ternyata kita dapat membuktikan bahwa himpunan itu tidak ada.

\begin{thm}[Paradoks Russell]\ollabel{thm:russells-paradox}
	Tidak ada himpunan $R = \Setabs{x}{x \notin x}$.
\end{thm}

\begin{proof}
Jika $R = \Setabs{x}{x \notin x}$ ada, maka
$R \in R$ jika dan hanya jika $R \notin R$, dan ini merupakan
kontradiksi.
\end{proof}

\begin{tagblock}{novice}
\begin{explain}
Mari kita telaah bukti ini dengan lebih perlahan. Jika $R$ ada, masuk
akal untuk menanyakan apakah $R \in R$. Andaikan memang $R \in R$.
Himpunan $R$ didefinisikan sebagai himpunan semua himpunan yang tidak
merupakan !!{element}s dari dirinya sendiri. Jadi, jika $R \in R$,
maka $R$ sendiri tidak memiliki sifat yang mendefinisikan~$R$. Namun,
hanya himpunan yang memiliki sifat tersebut yang berada di dalam~$R$.
Karena itu, $R$ tidak mungkin merupakan !!a{element} dari~$R$, yaitu
$R \notin R$. Akan tetapi, $R$ tidak mungkin sekaligus merupakan dan
bukan merupakan !!a{element} dari~$R$, sehingga kita memperoleh
kontradiksi.

Karena asumsi bahwa $R \in R$ menghasilkan kontradiksi, kita mendapat
$R \notin R$. Namun, hal ini juga menghasilkan kontradiksi!{} Jika
$R \notin R$, maka $R$ sendiri memiliki sifat yang mendefinisikan
$R$, sehingga $R$ akan menjadi !!a{element} dari $R$, sama seperti
semua himpunan lain yang tidak beranggotakan dirinya sendiri. Sekali
lagi, himpunan itu tidak mungkin sekaligus bukan merupakan dan merupakan
!!a{element} dari~$R$.
\end{explain}
\end{tagblock}

\begin{digress}
Bagaimana kita menyusun teori himpunan yang tidak terjerumus ke dalam
Paradoks Russell, yaitu teori yang tidak membuat klaim
\emph{inkonsisten} bahwa $R = \Setabs{x}{x \notin x}$ ada? Kita perlu
menetapkan aksioma-aksioma yang memberikan syarat sangat tepat untuk
menyatakan kapan himpunan ada (dan kapan himpunan tidak ada).
	
Teori himpunan yang digambarkan secara garis besar dalam bab ini tidak
melakukan hal tersebut. Teori ini benar-benar \emph{naif}. Teori ini
hanya menyatakan bahwa himpunan mematuhi ekstensionalitas dan bahwa,
jika kita mempunyai beberapa himpunan, kita dapat membentuk gabungan,
irisan, dan sebagainya. Teori himpunan dapat dikembangkan secara lebih
ketat daripada ini. \oliflabeldef{cumul:::part}{Pembahasan ketat itu
akan disediakan dalam Bagian \olref[cumul][][]{part}. Untuk saat ini,
kita akan melanjutkan secara naif dan dengan hati-hati berusaha
menghindari kontradiksi.}{}
\end{digress}


\end{document}
